\subsection{Existing Uses of Preemption in Dataplane OS}

State-of-the-art dataplane OS schedulers use preemptive scheduling to reduce the impact of head-of-line blocking on tail latency.
Preemption can appear on both control path and datapath.

\paragraph{Control path preemption.}
Scheduler balances load across cores or adjusts core allocation for applications. 
This aims to maintain low core utilization to indirectly avoid head-of-line blocking and reduce queuing delays.
This type of preemption have more relaxed timing requirements since it does not directly affect latency-critical tasks.

\paragraph{Data path preemption.}
Scheduler directly preempts running tasks on the latency-critical datapath.
This aims to approximate Processor Sharing.
This type of preemption requires precise timing to ensure latency-sensitive tasks receive timely execution and low overhead to avoid excessive disruption to the datapath.

Existing works explore various preemption mechanisms, including POSIX signals, Inter-Processor Interrupts (IPIs), and cache-coherent message passing.
\yihan{We compare the overhead of these mechanisms in table xx.}

\subsection{Scalability Dilemma of Centralized Preemption}

We now discuss the scalability challenges of using a centralized preemption dispatcher to preempt remote worker cores.

Centralized preemption dispatcher implementation spectrum has two extremes:
\begin{enumerate*}[label=(\roman*)]
  \item \textit{Per-core preemption}, where the dispatcher issues preemptions to individual cores immediately to ensure precise preemption, and
  \item \textit{Batch-all preemption}, where the dispatcher issues preemptions to all cores at fixed intervals to minimize preemption issuing overhead.
\end{enumerate*}
\yihan{Explain why intermediate designs are bounded by these two}

We use an analytical model to illustrate the scalability dilemma of centralized preemption dispatchers.

\paragraph{System setup}
We consider a system with $n$ worker cores and preemption quantum $q$.
We define \textit{preemption drift} $\delta_i$ as the difference between the ideal preemption time (preemption quantum $q$) and actual preemption time when the preemption is received by core $i$.
The worst-case preemption drift across all cores is defined as $\Delta = \max_{i=1}^{n} \delta_i$.

% \paragraph{Preemption precision and preemption overhead}
% Centralized dispatchers face a fundamental trade-off between preemption precision and preemption overhead.
% In the instant per-core preemption extreme, the dispatcher issues preemptions to each core immediately.
% This ensures high preemption precision of $O(1)$, but incurs preemption overhead of $O(n)$, since the dispatcher must issue preemptions to all $n$ cores individually.
% In contrast, in the fixed-interval batch-all preemption extreme, the dispatcher issues preemptions to all cores at fixed intervals.
% This minimizes preemption overhead to $O(1)$, but results in low preemption precision of $O(q)$, since a core may have to wait up to $q$ time before receiving a preemption.

% \paragraph{Preemption precision and core efficiency}

\subsection{Hardware Isolation}